% Phần riêng: Vấn đề 3 — Phân bố bộ nhớ (được \input từ báo cáo chính)
% Biên dịch độc lập: pdflatex van_de_3_phan_bo_bo_nho_standalone.tex

\section{Vấn đề 3: Phân bố bộ nhớ FLASH/RAM}
\label{sec:van_de_3_mem}

\subsection{Phát biểu vấn đề}

Triển khai SOPWM trên TMS320F280049C không chỉ cần ``có đủ byte'' --- mỗi loại dữ liệu phải nằm đúng \textbf{loại bộ nhớ vật lý} và đúng \textbf{section liên kết}, nếu không hệ thống sẽ lỗi im lặng hoặc build thất bại. Đây là bài toán cấu trúc, không phải tối ưu phụ.

\paragraph{(a) Khối lượng LUT góc lớn so với RAM LS.}
Mỗi cặp $(N, m)$ cần lưu trước $M=(N-1)/2$ góc chuyển mạch (độ). Với năm mức $N$ và $100$ bước $m$, tổng cộng:
\begin{equation}
  S_{\mathrm{LUT,naive}} = 100 \times 4 \times (3+4+5+6+7) = \SI{10000}{byte}.
  \label{eq:lut_total}
\end{equation}
Nếu đặt toàn bộ LUT vào RAM (như bảng sin/cos tham chiếu trong module SVPWM trước đó), section \texttt{.bss}/\texttt{.const} dễ \textbf{tràn RAMLS5} --- lỗi linker đã ghi nhận trong \texttt{log.md}.

\paragraph{(b) Bảng lịch runtime không được nằm FLASH.}
Sau khi tra LUT, \texttt{SOPWM\_BuildSchedule()} ghi liên tục vào bảng $\texttt{sopwm\_sched}$ --- nguồn burst của DMA. Dữ liệu này:
\begin{itemize}[noitemsep]
  \item thay đổi mỗi chu kỳ cơ bản (\SI{500}{\hertz});
  \item phải nằm trong vùng RAM mà \textbf{DMA bus master} truy cập được (GS RAM trên F28004x);
  \item cần \textbf{hai bản sao} (double buffer) để CPU ghi buffer inactive trong khi DMA đọc buffer active.
\end{itemize}
Đặt nhầm vào RAM LS hoặc FLASH sẽ khiến DMA đọc sai, hoặc ghi CMP ra giá trị \texttt{0xFFFF}.

\paragraph{(c) Xung đột nhu cầu truy cập.}
Cùng một lúc hệ thống cần:
\begin{center}
\begin{tabular}{@{}lll@{}}
  \toprule
  Thành phần & Truy cập & Ràng buộc vùng \\
  \midrule
  LUT góc & CPU đọc, không đổi & FLASH \texttt{.const} \\
  \texttt{sopwm\_sched} & CPU ghi + DMA đọc & RAMGS, section riêng \\
  Stack / biến debug & CPU R/W & RAMLS / \texttt{.bss} \\
  Mã \texttt{BuildSchedule} & CPU thực thi & FLASH \texttt{.text} \\
  \bottomrule
\end{tabular}
\end{center}

\paragraph{(d) Chi phí bộ nhớ tăng theo tham số thiết kế.}
Nếu lưu thẳng $4M$ góc đã mở rộng (thay vì $M$ góc Q1), dung lượng LUT nhân bốn:
\begin{equation}
  S_{\mathrm{LUT,full}} = 4 \times S_{\mathrm{LUT,naive}} = \SI{40}{kB}.
\end{equation}
Tương tự, bảng lịch tăng tuyến tính theo $N_{\mathrm{carr}}$ và số pha:
\begin{equation}
  S_{\mathrm{sched}} = N_{\phi} \times N_{\mathrm{buf}} \times N_{\mathrm{carr}} \times \sizeof(\texttt{SOPWM\_CycleCmp\_t}).
  \label{eq:sched_size}
\end{equation}
Với $N_{\phi}=3$, $N_{\mathrm{buf}}=2$, $N_{\mathrm{carr}}=100$, mỗi phần tử 4 byte:
\begin{equation}
  S_{\mathrm{sched}} = 3 \times 2 \times 100 \times 4 = \SI{2400}{byte}.
  \label{eq:sched_2400}
\end{equation}

\noindent\textbf{Tóm lại,} vấn đề cần giải quyết là: \emph{phân tách dữ liệu tĩnh (LUT) và dữ liệu động (bảng lịch DMA), map đúng section trong linker, đồng thời tối giểm footprint nhờ đối xứng quarter-wave --- không tính góc runtime trên CPU.}

\subsection{Cấu hình bộ nhớ trong dự án}

\subsubsection{Bản đồ vật lý F280049C (trích \texttt{28004x\_generic\_flash\_lnk.cmd})}

\begin{table}[htbp]
  \centering
  \caption{Vùng bộ nhớ on-chip liên quan SOPWM.}
  \label{tab:physical_mem}
  \begin{tabular}{@{}llll@{}}
    \toprule
    Symbol linker & Địa chỉ gốc & Dung lượng & Vai trò trong SOPWM \\
    \midrule
    \texttt{FLASH\_BANK0\_SEC4} & \texttt{0x084000} & \SI{4}{\kibi\byte} & Section \texttt{.const} --- LUT \\
    \texttt{RAMGS0} & \texttt{0x00C000} & \SI{8}{\kibi\byte} & Section \texttt{ramgs0} --- \texttt{sopwm\_sched} \\
    \texttt{RAMLS5} & \texttt{0x00A800} & \SI{2}{\kibi\byte} & \texttt{.bss}, \texttt{.data} --- biến toàn cục \\
    \texttt{RAMM1} & \texttt{0x000400} & $\approx$\SI{1}{\kibi\byte} & Stack \\
    \bottomrule
  \end{tabular}
\end{table}

\subsubsection{Cấu hình section trong file liên kết}

File \texttt{28004x\_generic\_flash\_lnk.cmd} khai báo hai section \textbf{do dự án thêm} (ngoài mặc định C2000Ware):

\begin{lstlisting}[language=]
/* PAGE 1 -- MEMORY block */
RAMGS0 : origin = 0x00C000, length = 0x002000   /* 8 KiB */

/* SECTIONS */
.const   : > FLASH_BANK0_SEC4,  PAGE = 0, ALIGN(4)   /* const float LUT */
ramgs0    : > RAMGS0,            PAGE = 1             /* DMA source table */
ramgs1    : > RAMGS1,            PAGE = 1             /* du phong mo rong */
\end{lstlisting}

\paragraph{Lý do chọn \texttt{FLASH\_BANK0\_SEC4} cho \texttt{.const}.}
Sector \SI{4}{\kibi\byte} đủ chứa $\approx$\SI{10}{\kibi\byte} LUT (chiếm $\approx 24\%$ sector), tách biệt khỏi \texttt{.text} (SEC2/3/5) và \texttt{.cinit} (SEC1), tránh phải gom mã và hằng số vào cùng sector gây khó mở rộng.

\paragraph{Lý do chọn \texttt{RAMGS0} cho bảng lịch.}
Theo TRM F28004x, DMA có thể truy cập Global Shared RAM (GS). \texttt{sopwm\_sched} chiếm \SI{2400}{byte} trên \SI{8192}{byte} RAMGS0 ($\approx 29\%$), còn dư cho biến DMA khác. Section \texttt{ramgs1} được khai báo sẵn trên RAMGS1 (\SI{8}{\kibi\byte}) cho mở rộng sau này.

\subsubsection{Ràng buộc build CCS}

\begin{itemize}[noitemsep]
  \item Cấu hình \textbf{CPU1\_FLASH} (hoặc tương đương) dùng \texttt{28004x\_generic\_flash\_lnk.cmd}.
  \item \texttt{Flash\_initModule(..., DEVICE\_FLASH\_WAITSTATES)} trong \texttt{device.c}: \texttt{WAITSTATES = 4} @ \SI{100}{\mega\hertz} --- LUT đọc từ FLASH vẫn đủ nhanh vì \texttt{BuildSchedule} chỉ chạy \SI{500}{\hertz}, không phải \SI{50}{\kilo\hertz}.
  \item DMA \texttt{srcWrapSize = 200} trong SysConfig khớp $2 \times N_{\mathrm{carr}}$ word nguồn (double buffer liên tiếp trong RAM).
\end{itemize}

\subsection{Đóng góp của mã nguồn}

Bảng~\ref{tab:code_contrib} tóm tắt cách code giải quyết vấn đề phân bố bộ nhớ.

\begin{table}[htbp]
  \centering
  \caption{Đóng góp mã nguồn cho bài toán bộ nhớ.}
  \label{tab:code_contrib}
  \small
  \begin{tabular}{@{}p{3.2cm}p{4.2cm}p{5.8cm}@{}}
    \toprule
    Thành phần & File & Đóng góp \\
    \midrule
    LUT tĩnh &
    \texttt{sopwm.c} &
    Năm mảng \texttt{const float SOPWM\_LUT\_N*} ($\approx$\SI{10}{kB}) tự động vào \texttt{.const}/FLASH; không dùng malloc, không copy sang RAM lúc khởi động. \\
    \addlinespace
    Tra cứu $O(1)$ &
    \texttt{sopwm.h} &
    \texttt{SOPWM\_GetAngles()} \texttt{static inline}: map $m \rightarrow i_m=\mathrm{round}(100m-1)$, \texttt{switch(N)} --- thời gian xác định, không stack frame. \\
    \addlinespace
    Placement DMA &
    \texttt{sopwm.c} &
    \texttt{\#pragma DATA\_SECTION(sopwm\_sched, "ramgs0")} buộc linker đặt bảng vào RAMGS0; comment ghi rõ yêu cầu linker. \\
    \addlinespace
    Kiểu dữ liệu gọn &
    \texttt{sopwm.h} &
    \texttt{SOPWM\_CycleCmp\_t} = 2$\times$\texttt{uint16\_t} (4 byte/slot), không dùng \texttt{float} trong bảng DMA --- giảm $S_{\mathrm{sched}}$ và khớp burst DMA (2 word). \\
    \addlinespace
    Double buffer &
    \texttt{sopwm.c/h} &
    \texttt{sopwm\_sched[3][2][SOPWM\_N\_CARR]}; \texttt{sopwm\_sched\_active}, \texttt{sopwm\_sched\_next}; \texttt{BuildSchedule} ghi \texttt{[next]}, ISR hoán đổi --- tránh ghi đè vùng DMA đang đọc. \\
    \addlinespace
    Tiết kiệm FLASH &
    \texttt{sopwm.c} &
    \texttt{SOPWM\_BuildSchedule()}: chỉ đọc $M$ góc Q1 từ LUT, mở rộng $4M$ góc bằng~\eqref{eq:quarter_expand} trên stack (\texttt{all\_angles[28]}) --- không lưu $4M$ góc/cấp $m$ trong FLASH. \\
    \addlinespace
    Linker &
    \texttt{28004x\_generic\_flash\_lnk.cmd} &
    Map \texttt{.const} sang FLASH, \texttt{ramgs0} sang RAMGS0; tách LUT khỏi RAMLS tránh tràn như SVPWM. \\
    \bottomrule
  \end{tabular}
\end{table}

\subsubsection{Minh hoạ placement trong \texttt{sopwm.c}}

LUT --- compiler đặt vào \texttt{.const} (FLASH):
\begin{lstlisting}[language=C]
const float SOPWM_LUT_N7[SOPWM_M_COUNT][3] = {
    /* m=0.01 */ {60.118f, 80.177f, 80.364f},
    /* ... 100 rows ... */
};
\end{lstlisting}

Bảng lịch --- placement explicit sang RAMGS:
\begin{lstlisting}[language=C]
#pragma DATA_SECTION(sopwm_sched, "ramgs0")
SOPWM_CycleCmp_t sopwm_sched[3][2][SOPWM_N_CARR];
\end{lstlisting}

BuildSchedule --- đọc FLASH, ghi RAMGS (buffer inactive):
\begin{lstlisting}[language=C]
void SOPWM_BuildSchedule(float m, uint16_t N, uint16_t tbprd)
{
    float base_deg[7];
    float all_angles[28];   /* 4*7 max -- stack, khong FLASH */
    n_base = SOPWM_GetAngles(m, N, base_deg);  /* doc LUT FLASH */
    /* ... mo rong quarter-wave vao all_angles[] ... */
    sopwm_bin_lut(all_angles, total, tbprd,
                  sopwm_sched[0][sopwm_sched_next]);  /* ghi RAMGS */
}
\end{lstlisting}

\subsection{Sơ đồ phân vùng logic}

\begin{figure}[htbp]
  \centering
  \begin{tikzpicture}[
    scale=0.92,
    every node/.style={font=\small},
    flash/.style={draw, fill=orange!15, minimum width=4.8cm, minimum height=1.1cm, align=center},
    ramgs/.style={draw, fill=green!15, minimum width=4.8cm, minimum height=1.1cm, align=center},
    raml/.style={draw, fill=blue!10, minimum width=4.8cm, minimum height=0.9cm, align=center},
    arr/.style={-{Stealth}, thick}
  ]
    \node[flash] (fconst) at (0,2.4) {\textbf{FLASH SEC4}\\ \texttt{.const}: \texttt{SOPWM\_LUT\_N7..N15}\\ $\approx$\SI{10}{\kibi\byte}};
    \node[flash] (ftext) at (0,1.0) {\textbf{FLASH SEC2/3/5}\\ \texttt{.text}: \texttt{BuildSchedule}, ISR};

    \node[ramgs] (gs) at (7,2.4) {\textbf{RAMGS0} (\texttt{ramgs0})\\ \texttt{sopwm\_sched[3][2][100]}\\ \SI{2400}{byte} --- nguon DMA};
    \node[raml] (ls) at (7,1.0) {\textbf{RAMLS5}\\ \texttt{.bss}/\texttt{.data}: \texttt{sopwm\_m\_cmd}, debug};

    \draw[arr] (fconst) -- node[above, align=center, font=\scriptsize] {CPU doc\\@ 500 Hz} (gs);
    \draw[arr] (ftext) -- node[above, align=center, font=\scriptsize] {CPU ghi\\buffer next} (gs);
    \draw[arr, dashed] (gs.east) -- ++(1.1,0) node[right, align=left, font=\scriptsize] {DMA burst\\50 kHz};
  \end{tikzpicture}
  \caption{Phân vùng logic: LUT tĩnh trên FLASH; bảng lịch động trên RAMGS; CPU là cầu nối, DMA đọc trực tiếp RAMGS.}
  \label{fig:mem_map}
\end{figure}

\subsection{Chi tiết dung lượng LUT}

\begin{table}[htbp]
  \centering
  \caption{Dung lượng từng bảng \texttt{SOPWM\_LUT\_N*}.}
  \label{tab:lut_breakdown}
  \begin{tabular}{@{}lccr@{}}
    \toprule
    Bảng & Kích thước & Phần tử & Bytes \\
    \midrule
    \texttt{SOPWM\_LUT\_N7}  & $100\times 3$ & 300  & 1200 \\
    \texttt{SOPWM\_LUT\_N9}  & $100\times 4$ & 400  & 1600 \\
    \texttt{SOPWM\_LUT\_N11} & $100\times 5$ & 500  & 2000 \\
    \texttt{SOPWM\_LUT\_N13} & $100\times 6$ & 600  & 2400 \\
    \texttt{SOPWM\_LUT\_N15} & $100\times 7$ & 700  & 2800 \\
    \midrule
    \textbf{Tổng} & & 2500 float & \textbf{10000} \\
    \bottomrule
  \end{tabular}
\end{table}

Chỉ số tra cứu (code):
\begin{equation}
  i_m = \mathrm{round}(100\,m - 1), \qquad i_m \in [0,\,99].
  \label{eq:lut_index_v3}
\end{equation}

\subsection{Kết quả và hệ quả thiết kế}

\begin{enumerate}
  \item \textbf{Linker thành công} với LUT trên FLASH và bảng lịch trên RAMGS --- không tràn RAMLS như phiên bản SVPWM.
  \item \textbf{DMA hoạt động ổn định} vì nguồn nằm đúng GS RAM; \texttt{srcWrapSize=200} khớp layout $[3\,\text{pha}][2\,\text{buf}][100]$ khi mỗi kênh DMA phục vụ một pha.
  \item \textbf{Footprint LUT giảm $\times 4$} nhờ lưu $M$ góc Q1 + mở rộng runtime; footprint bảng lịch giảm nhờ \texttt{uint16\_t} thay vì \texttt{float} cho CMP.
  \item \textbf{Chi phí mở rộng có thể dự báo:} tăng $N_{\mathrm{carr}}$ từ 100 lên 200 sẽ đẩy $S_{\mathrm{sched}}$ lên \SI{4800}{byte}; tăng số mức $m$ từ 100 lên 1000 sẽ nhân mười dung lượng FLASH LUT.
\end{enumerate}

\noindent\textit{Kiểm tra sau build:} trong CCS, cửa sổ \texttt{Memory Browser} xác nhận địa chỉ \texttt{sopwm\_sched} thuộc vùng \texttt{0x00C000}--\texttt{0x00DFFF}; map file liệt kê \texttt{SOPWM\_LUT\_N*} trong section \texttt{.const}.
